% Standalone source-record documentation; no research attempts.
\documentclass[11pt]{article}
\usepackage[margin=25mm]{geometry}
\usepackage{fontspec,unicode-math}
\setmainfont{Latin Modern Roman}
\setmathfont{Latin Modern Math}
\usepackage{amsmath,amssymb,mathtools,xurl,hyperref}
\hypersetup{colorlinks=true,urlcolor=blue,linkcolor=blue}
\setcounter{secnumdepth}{0}
\setlength{\parindent}{0pt}
\setlength{\parskip}{5pt}
\setlength{\emergencystretch}{3em}
\newcommand{\sref}[2]{\hyperlink{source:#1:#2}{[S#2]}}
\newcommand{\cataloguescope}[1]{\paragraph{Source record.}#1}
\newcommand{\floor}[1]{\left\lfloor#1\right\rfloor}
\begin{document}
% UnsolvedMath ID 20001166; source catalogue code AIM-COMPUTATION-0004
\setcounter{section}{20001165}
\section{Exact ridge reduction for Lévy PIDEs}\label{20001166}
{\small\sffamily UnsolvedMath ID 20001166 · Computer Science}
\subsection{Definitions and mathematical statement}

\paragraph{Shared notation.}Standard number systems are $\mathbb N,\mathbb Z,\mathbb Q,\mathbb R,\mathbb C$. Specialized notation and index ranges follow the cited source and the definitions below.


A partial integro-differential equation combines derivatives of an unknown function with a nonlocal integral term. A Lévy measure $\nu$ describes jump sizes and satisfies $\int(1\wedge\|z\|^2)\,\nu(dz)<\infty$. A viscosity subsolution or supersolution satisfies the corresponding differential inequalities when touched by suitable test functions. Comparison bounds every subsolution by every supersolution with compatible boundary data, often yielding uniqueness.
\paragraph{Record scope and qualifications.}
The catalogue title is retained as metadata. Named certificates, contracts, bounds or frameworks in that title are not asserted as results of the original source extract.
This is a participant research agenda rather than a fixed PIDE. Dynamics, Lévy measures, boundary conditions, and the requested numerical or regularity guarantee are not specified. The supplied catalogue marks this record open; its present research status was not independently checked.
The source text below is rendered literally where necessary; damaged formulas and composite extracts are not silently repaired into a new mathematical claim.
\paragraph{Statement recorded in the supplied source.}
A.1 Arisawa, Mariko

We would like to introduce a viscosity solution's approach to a class of second-order partial-integro-differential equations (PIDE, in short) arising in the mathematical finance. The usual brownian models for the processes of the asset prices are replaced by the jump-diffusion processes. The derived PIDEs contain in the integral term, various Le'vy measures which correspond to the frequencies and the sizes of the jumps. In addition, according to the financial settings such as the American option, the look back option, the problem with transaction costs, a variety of nonlinear PIDEs are engendered. We shall introduce a new definition of the viscosity solution for such problems, and shall give the comparison, the existence, and the regularity of the solutions of viscosity solutions in this framework. As far as the high-dimensional numerical analysis, a very few is known at this stage, while from the practical view point a lot of important problems should be explored, such as the modelization of the high-dimensional Le'vy measures arising in the mathematical finances. This is the problem which I would like to discuss with the participants of the workshop. \sref{20001166}{2}
\paragraph{Source status.}
The supplied catalogue labels this source record open; that metadata was checked against the supplied snapshot on 2026-10-11. The label is preserved for provenance and is not an independent verification of present open status or mathematical correctness.
\subsection{Short English statement}
The participant proposes viscosity-solution theory and high-dimensional numerical methods for jump models in finance.
\subsection{Sources}
\begin{itemize}
\item[\hypertarget{source:20001166:1}{S1}] Ulam AI and the UnsolvedMath Contributors, UnsolvedMath: A Curated Collection of Open Mathematics Problems, record 20001166 (AIM-COMPUTATION-0004). Catalogue curation and metadata: CC BY 4.0. The source extract is retained; definitions, English summaries and qualifications are editorial additions.
\url{https://huggingface.co/datasets/ulamai/UnsolvedMath}.
\item[\hypertarget{source:20001166:2}{S2}] AIM Workshop Problem Lists, source linked by catalogue record AIM-COMPUTATION-0004.
\url{https://aimath.org/WWN/optimalcontrol/optimalcontrol.pdf}.
\end{itemize}
\paragraph{Preparation and review.}
The original source and UnsolvedMath contributors retain attribution for the source extract. Codex (GPT-6; exact serving revision not exposed) assisted with definitions, English summaries, scope notes and formatting. No new human mathematical authorship or independently verified solution is asserted. Permanent source ID 20001166 and code AIM-COMPUTATION-0004 checked against the complete supplied registry and current repository. Distinct upstream records, including repeated questions and broad or damaged extracts, are retained for complete coverage of the source CS category. A repeated record is not a new independent conjecture. Complete mathematical-equivalence and current-literature review has not been performed. Preparation and final review follow contribution prompts C and D, with source gaps disclosed. Structural and rendering checks do not establish a new theorem or open status.
\label{end:20001166}
\end{document}
